\section{Rapport synthétique}

\subsection{Cas clinique}
Patient de $79$ ans est hospitalisé pour décompensation cardiaque avec fraction d’éjection conservée à $66$\%, Heart Failure with preserved Ejection Fraction (HFpEF) à la suite de l’administration de la première dose d’Azacitidine (Vidaza). Azacitidine est un agent déméthylant de l’ADN prescrit dans le cadre d’une leucémie myéloïde aiguë.

Antécédents médicaux pertinents

\begin{itemize}[label=\textbullet]
    \item Diabète de type $2$ ;
    \item Syndrome métabolique ;
    \item Cardiopathie hypertensive avec dysfonction diastolique \todoH[FloraStyle][red!30]{\ !!! faut savoir c'est quoi la classification !}{\mbox{franche de grade $2$}} ;
    \item Dermo-hypodermite au niveau prétibial gauche ;
    \item Fibrillation atriale ;
    \item Infarctus du myocarde 
    \item BPCO GOLD II
    \item Leucémie myéloïde aiguë
    \end{itemize}

Antécédents chirurgicaux pertinents

\begin{itemize}[label=\textbullet]
    \item Mise en place d’un Pacemaker DDD ;
    \end{itemize}

Allergies et assuétudes

\begin{itemize}[label=\textbullet]
    \item Pas d’allergie connue à ce jour ;
    \item Tabac : $50$ UPA ($2$ paquets/j pendant $25$ ans) arrêté depuis une vingtaine d’années ;
    \item Alcool : occasionnelle ;
    \item Il nie toute autre consommation ;
    \end{itemize}


Médication (en septembre $2026$) reprise dans le~\cref{tab:medication_patient}.

% ====================================================================================
% VERSION 3 : PAYSAGE / PENCHÉ (STYLE LES_HYPOTHESES_DIAGNOSTIQUES)
% Page pivotée automatiquement dans le visualiseur PDF avec grille complète (| et \hline)
% ====================================================================================
\begin{landscape}
\begingroup
\footnotesize
\setlength{\tabcolsep}{5pt}
\renewcommand{\arraystretch}{1.05}

\begin{xltabular}{\linewidth}{|>{\raggedright\arraybackslash\bfseries}p{4.2cm}|>{\raggedright\arraybackslash}p{3.6cm}|>{\raggedright\arraybackslash}p{5.8cm}|>{\raggedright\arraybackslash}X|}
    \caption{Médication du patient (septembre 2026)} \label{tab:medication_patient} \\
    \hline
    \rowcolor{gray!15}
    \textbf{DCI} & \textbf{Nom commercial \& Posologie} & \textbf{Classe pharmacothérapeutique} & \textbf{Indication clinique} \\
    \hline
    \endfirsthead

    \caption*{Tableau~\ref{tab:medication_patient} (suite) -- Médication du patient (septembre 2026)} \\
    \hline
    \rowcolor{gray!15}
    \textbf{DCI} & \textbf{Nom commercial \& Posologie} & \textbf{Classe pharmacothérapeutique} & \textbf{Indication clinique} \\
    \hline
    \endhead

    \hline
    \multicolumn{4}{|r|}{\footnotesize\textit{Suite du tableau à la page suivante...}} \\
    \hline
    \endfoot

    \hline
    \endlastfoot

    Ézétimibe / Atorvastatine & 
    Atozet \qty{10}{\mg} / \qty{40}{\mg} \newline 1 cp/j & 
    Hypolipémiant combiné (inhibiteur sélectif de l'absorption du cholestérol + statine de haute intensité) & 
    Prévention cardiovasculaire secondaire. \\
    \hline

    Bumétanide & 
    Burinex \qty{1}{\mg} \newline 1 cp/j & 
    Diurétique de l'anse de Henle & 
    Traitement déplétif de la surcharge hydrosodée : contrôle des pressions de remplissage ventriculaire gauche élevées liées à la cardiopathie hypertensive avec dysfonction diastolique de grade 2. \\
    \hline

    Bisoprolol & 
    Bisoprolol \qty{10}{\mg} \newline ½ cp/j (\qty{5}{\mg}/j) & 
    Bêtabloquant $\beta_1$ cardiosélectif & 
    Cardiopathie hypertensive \& contrôle hémodynamique : baisse des résistances vasculaires périphériques, allongement du temps de remplissage diastolique par ralentissement de la fréquence cardiaque, et freinage de l'hyperactivité sympathique. \\
    \hline

    Amiodarone & 
    Cordarone \qty{200}{\mg} \newline 1 cp/j & 
    Antiarythmique de classe III de Vaughan-Williams & 
    Prévention des récidives de tachyarythmie supraventriculaire. \\
    \hline

    Fluticasone furoate / Uméclidinium / Vilantérol & 
    Trelegy Ellipta 92/55/22\,\unit{\ug} \newline 1 inh/j & 
    Trithérapie inhalée fixe : Corticoïde inhalé (CSI) + Anticholinergique à longue durée d'action (LAMA) + $\beta_2$-mimétique à longue durée d'action (LABA) & 
    Traitement de fond de la BPCO. \\
    \hline

    Bromure d'ipratropium / Fénotérol & 
    Duovent \newline (posologie à la demande) & 
    Bronchodilatateur d'action rapide combiné : Anticholinergique à courte durée d'action (SAMA) + $\beta_2$-mimétique à courte durée d'action (SABA) & 
    Traitement de secours / à la demande : levée immédiate du bronchospasme et prise en charge des accès dyspnéiques paroxystiques dans la BPCO. \\
    \hline

    Apixaban & 
    Eliquis \qty{5}{\mg} \newline 1 cp 2$\times$/j & 
    Anticoagulant oral direct (AOD) & 
    Thromboprophylaxie systémique liée à l'arythmie supraventriculaire associée (score $\mathrm{CHA_2DS_2}\text{-}\mathrm{VA}$ élevé). \\
    \hline

    Cholécalciférol (Vitamine $\mathrm{D_3}$) & 
    Vista $\mathrm{D_3}$ \newline 1 cp/j & 
    Vitamine D / régulateur du métabolisme phosphocalcique & 
    Prévention de l'ostéopénie / ostéoporose. \\
    \hline

    Zolpidem & 
    Stilnox \qty{10}{\mg} \newline 1 cp le soir & 
    Hypnotique apparenté aux benzodiazépines & 
    Traitement symptomatique de courte durée de l'insomnie d'endormissement. \\
    \hline

    Azacitidine & 
    Vidaza \newline Cure palliative par voie sous-cutanée : 7 jours/mois pendant 28 jours & 
    Agent déméthylant de l'ADN & 
    Traitement de la leucémie myéloïde aiguë non éligible à une chimiothérapie d'induction intensive. \\
    \hline

    Dapagliflozine & 
    Forxiga \qty{10}{\mg} \newline 1 cp/j & 
    Inhibiteur du cotransporteur sodium-glucose de type 2 (iSGLT2) & 
    Triple indication : \newline
    \textbullet~Diabète de type 2 ; \newline
    \textbullet~Cardioprotection ; \newline
    \textbullet~Néphroprotection. \\
    \hline

    Oxygénothérapie chronique & 
    \qty{2}{\L\per\minute} par oxyconcentrateur & 
    Support ventilatoire & 
    Support pour l'hypoxémie nocturne et les troubles de diffusion dans le cadre de la BPCO. \\
\end{xltabular}
\endgroup
\end{landscape}


Contexte psyho social : 
Le patient vit au domicile avec son épouse. Elle est très présente pour le patient. Le patient a $2$ enfants et $4$ petits-enfants. 
Le patient se sent très limité, dyspnéique et trouve le temps long en chambre.  Une thymie triste domine dans le discours du patient. 

Histoire clinique :

Le patient est hospitalisé dans le service d’onco-hématologie pour insuffisance respiratoire de type I sur décompensation cardiaque. Suite à l’administration de la première dose d’Azacitidine (Vidaza) dans le contexte de sa leucémie myéloïde aiguë, le patient a présenté une dyspnée de grade IV NYHA évolutive avec sensation d’oppression thoracique. Ceci s’ajoutant à des signes cliniques caractéristiques de la décompensation cardiaque décrits dans l’examen clinique. Il a été décrit dans la littérature une augmentation significative d’événement cardio-pulmonaire lorsque le patient a des polymorbidités de ces systèmes. De plus, la dyspnée est répertoriée comme un effet indésirable très fréquent selon \textcite{ema_azacitidine_accord}.

Paramètres d’entrée :
\begin{itemize}[label=\textbullet]
    \item Tension artérielle couché : $96/70$
    \item Saturation : $91$\% sous $4$L/min d’oxygène
    \item Température : $36,8$°C
    \item Pouls : $69$/min
    \item Taille : $167$cm
    \item Poids : $58$kg
    \item BMI : $20,8$ kg/m$^2$
    \end{itemize}

Examen clinique d’entrée :
    \begin{itemize}[label=\textbullet]
    \item Apparence : pâle.
    \item Examen neurologique : le patient est orienté dans l’espace et le temps. Score de Glasgow de $15/15$.
    \item Examen pulmonaire : eupnéique sous 4L/min d’oxygène. A l’auscultation, murmure vésiculaire diffus et symétrique sans sibilance. Crépitants bibasilaires entendus.
    \item Examen cardiaque et vasculaire : les jugulaires sont turgescentes avec un reflux hépato-jugulaire. Présence d’œdème des membres inférieurs de façon bilatérale avec prise du Godet à la palpation. A l’auscultation, B$1$B$2$ réguliers. Pas de souffle ni galop.
    \item Examen abdominal : péristaltisme normal à l’auscultation. Abdomen souple, dépressible, légèrement sensible en fosse iliaque gauche à la palpation. Pas d’organomégalie palpée.
    \end{itemize}


\subsection{Prise en charge hospitalière}
Volet cardio-pneumologique : 
\begin{itemize}[label=\textbullet]
    \item \underLine{A l’admission} : présentation clinique d’une décompensation cardiaque claire. Notamment des signes de décompensation cardiaque droite avec des œdèmes des membres inférieurs prenant le Godet, des jugulaires saillantes et un reflux hépato-jugulaire ainsi que des signes d’une décompensation cardiaque gauche avec des crépitants bibasilaires à l’auscultation. De plus, la biologie appuie cette étiologie avec un dosage de NT Pro-BNP à $2055.0$ pg/mL et la radiographie de thorax a démontré une cardiomégalie modérée avec surcharge vasculaire hilaire gauche.
    \item  \underLine{Avis cardiologique requis} : suite à la réalisation d’une échographie cardiaque, nous pouvons conclure que la fonction ventriculaire gauche est conservée et la dynamique régionale normale. Absence, donc, d’argument en faveur d’une élévation des pressions de remplissage. Par contre, une insuffisance tricuspidienne modérée et une hypertension artérielle pulmonaire est démontrée. Le cardiologue préconise une administration de Lasix IV et augmentation des doses du  Bumétanide  à $2,5$mg en conservant le Dapagliflozine. Une augmentation de l’oxygénothérapie à $4$L/min est proposée, à adapter selon la clinique et les paramètres du patient au cours de l’hospitalisation.
    \item \underLine{ Épisode de désaturation à $80$\% sous $3$L/min d’oxygène $4$ jours après l’admission} : augmentation de l’oxygénothérapie à $5$L/min a permis une saturation à $87$\%. Dans les heures qui ont suivi, le débit a pu redescendre à $3,5$L/min. Un nouveau bilan est débuté comprenant une radiographie du thorax, une biologie et un ECG. Les résultats principaux sont une majoration de l’infiltrat interstitiel bilatéral sur lésions de broncho-emphysème diffus sans foyer parenchymateux systématisé, une biologie avec D-Dimère à $750$ ng/mL, GOT $145$ U/L, proBNP $13 000$ pg/mL, Troponine $0,025$ µg/L à interpréter dans le contexte global de décompensation cardiaque, à savoir qu’il n’y avait aucune douleur thoracique rapportée. Une aérosolthérapie par Duovent a été débutée ainsi qu’une majoration du Bumétanide d’$1$mg et une ampoule de $0,5$mg de morphine qui a soulagé la symptomatologie. Bien que l'avis cardiologique ait suspecté une cause respiratoire prédominante (surinfection à Pseudomonas), l'ascension majeure du NT-proBNP à $13 000$ pg/mL et la réponse à la majoration du Burinex signent un tableau mixte. 
\end{itemize}

Volet infectiologique : 
\begin{itemize}[label=\textbullet]
    \item \underLine{Au cours de l’hospitalisation} : Monsieur présentait des expectorations purulentes avec péjoration de la dyspnée et une altération de l’état général. Toutefois, le patient est resté apyrétique durant l’ensemble de l’hospitalisation.
    \item \underLine{Avis infectiologique} : une surinfection pulmonaire a été suspectée. Il a été décidé de couvrir le patient par amoxicilline / acide clavulanique selon une antibiothérapie probabiliste. Cependant, à la suite de la positivité à Pseudomonas aeruginosa, bacille pyocyanique, la ceftazidime a été initiée pendant 7 jours.
    \end{itemize}

Volet Hémato-oncologique :
\begin{itemize}[label=\textbullet]
    \item \underLine{A l’admission} : interruption du traitement d’Azacitidine suite au contexte cardio-pulmonaire. 
    \item \underLine{Lors de l’hospitalisation} : surveillance étroite des marqueurs hématologiques. Une neutropénie débutante secondaire à Azacitidine a été observée.
    \item \underLine{A la sortie} : planification des soins et la poursuite du traitement palliatif de la leucémie myéloïde aiguë.
    \end{itemize}


Les équipes pluridisciplinaires intervenant au cours de l’hospitalisation de Monsieur :
\begin{itemize}[label=\textbullet]
    \item \underLine{L’équipe infirmière/aide-soignant} au cœur des soins. Le rôle majeur est dans la surveillance constante du patient et la transmission d’informations pertinentes à l’équipe médicale. Les infirmiers ont aussi l’opportunité de stimuler le patient à se mettre au fauteuil et apporter l’aide quotidienne nécessaire au patient. 
    \item \underLine{La diététicienne}  est intervenue dans ce cas, suite aux signes de fausses déglutitions rapportées par le patient et l’équipe infirmière. Les textures alimentaires ont pu être adaptées afin de correspondre aux besoins de Monsieur.
    \item \underLine{L’équipe des soins continus/palliatif et l’équipe sociale} ont collaboré ensemble afin de conclure à une solution envisageable par la patient et ses proches pour la suite de la prise en charge. Les conditions du projet thérapeutique ont également été rediscutées lors de l’hospitalisation afin de préserver la collaboration avec Monsieur et s’assurer une qualité de soins. 
    \item \underLine{Le psychologue de liaison} n’a pas été sollicité par refus du patient.
    \end{itemize}
    
Voir détails des examens complémentaires en \cref{annexe:examens_complementaires} en \cpageref{annexe:examens_complementaires}.

\subsection{Complexités rencontrées}
\begin{itemize}[label=\textbullet]
    \item \underLine{Complexité diagnostique} : déterminer l’étiologie de la dyspnée parmi un tableau clinique complexe. La multitude de pathologies que le patient cumule sont réparties dans les différents systèmes. Devoir agir dans l’urgence et décider un traitement est complexe.
    \item \underLine{Complexité thérapeutique} : aborder un patient avec autant de traitements concomitants est difficile et requiert une précaution supérieure face aux interactions médicamenteuses possibles. Savoir élaguer la liste médicamenteuse afin de garantir le plus de confort au patient est une tâche complexe.
    \item \underLine{Complexité psycho-sociale} : la décision prise par le patient ne correspond pas forcément à la solution idéale au niveau médical. Il faut donc apprendre à accepter l’autre et passer au-dessus de ce sentiment de frustration qui peut survenir.
    \end{itemize}
\subsection{Éclairage d’un aspect thérapeutique}
\subsubsection{Méthode}
\underLine{Formulation de la question clinique selon la méthode PICO :}
\begin{itemize}[label=\textbullet]
    \item POPULATION : patient polypathologique âgé atteint d'insuffisance cardiaque à fraction d'éjection préservée (HFpEF, FE $66$ \%) et diabète de type $2$.
    \item INTERVENTION : administration d'un inhibiteur du SGLT$2$.
    \item COMPARATEUR : traitement conventionnel sans iSGLT$2$ / placebo.
    \item OUTCOME : réduction des réhospitalisations pour décompensation cardiaque, mortalité et amélioration de la qualité de vie.
    \end{itemize}
\underLine{Base de données consultée :} PubMed.\\
\underLine{Equation de recherche utilisée :} ("Sodium-Glucose Transporter $2$ Inhibitors"[Mesh] OR "dapaglifloxzin") AND ("HFpEF"OR "preserved ejection fraction") AND ("Heart Failure"[Mesh] ).\\
\underLine{Critères de sélection :} date de publication <$5$ ans, texte intégral, revue systématique, méta-analyse.\\
\underLine{Date de consultation :} $25$ octobre $2026$.

\subsubsection{Résultat}
J’ai choisi la méta-analyse \textcite{minisy2025role}. Elle regroupe $9$ études randomisées entre $2015$ et $2025$ avec un ensemble de $20 000$ patients. La conclusion confirme l‘efficacité des inhibiteurs SGLT$2$ dans la prise en charge de l’insuffisance cardiaque avec fraction d’éjection préservée (HFpEF). La réduction du risque relatif de mortalité cardiovasculaire et le nombre d’hospitalisation pour insuffisance cardiaque serait de $17$\%. Ce résultat est obtenu de façon constante et significative à travers l’ensemble des études. Leur action sur la mortalité et la qualité de vie global tend à être favorable sans toutefois pouvoir être objectivée de façon significative. 

\subsubsection{Discussion}
Cette méta-analyse démontre clairement les bénéfices des inhibiteurs SGLT$2$ et donc justifie leur place au sein de la prise en charge de l’HFpEF. Plusieurs outils de qualités utilisés sont reconnus et validés tel que Cochrane RoB $2.0$, PRISMA $2020$ guidelines, évaluation GRADE, … ce qui est un avantage pour la méthode de l’analyse.

Toutefois, parmi les limites admises par les auteurs, nous retenons l’absence d’analyse en sous-groupe notamment pour dissocier les polymorbidités ou encore pour dissocier les différentes molécules utilisées. Ils l’expliquent par le choix de travailler sur des données agrégées et non individuelles. La définition stricte de la fraction d’éjection préservée n’a pas été respectée. Ils ont admis un plus faible pourcentage de fraction d’éjection allant jusqu’à $40$\% or en théorie, elle se limite à $50$\% ou plus. Au final, presque un tier de la cohorte a une fraction d’éjection <$50$\%. Le dernier point qu’ils mettent en lumière est l’hétérogénéité modérée au vu de l’ensemble des critères d’inclusion repris dans les différentes RCTs et les biais retrouvés dans certaines études notamment dû à un modèle non aveugle et aux données manquantes.

En analysant en détails la méta-analyse, j’ai pu remarquer d’autres défauts dans la conception de cette méta-analyse. D’une part, dans la sélection des RCTs, ils ont repris une étude, « DEFINE-HF » qui comprend dans son titre même « …...in patients with heart failure with reduced ejection fraction: DEFINE-HF trial » et certaines études, EMPEROR-Pooled et EMPEROR-Preserved, repris comprennent les mêmes patients et donc il s’agit d’un double comptage des patients ce qui fausse bien sûr les chiffres et l’importance des résultats obtenus. Un autre point primordial à aborder dans cette méta-analyse, spécifiquement lorsque la population est une population âgée, polypathologique, sont les effets indésirables et graves des inhibiteurs SGLT$2$. Or, cela n’a pas été analysé ou évoqué dans la méta-analyse.

En conclusion, il s’agit d’une méta-analyse ayant une base de données solide et une méthodologie de référence avec des outils fiables. Toutefois, plusieurs limites nous rappellent qu’il faut être vigilant quant à l’interprétation des résultats. Cependant, vu la tendance favorable des inhibiteurs SGLT$2$ sur le tableau clinique de l’insuffisance cardiaque à fraction préservée, il semble confirmer son rôle dans les guidelines. 

\subsubsection{Contextualisation} 
La décision médicale d’initier la bithérapie d’un diurétique de l’anse avec du Dapagliflozin dans le contexte de décompensation cardiaque était justifiée. De plus, cette méta-analyse appuie l’avis cardiologique donné à Monsieur, surtout vu la pertinence des outcomes tel que la réduction d’hospitalisation pour Monsieur désireux d’un projet de soins palliatifs à domicile.

Bien sûr, une surveillance clinique était requise compte tenu du nombre accru de médicaments et pathologies dont le patient souffrait afin de prévenir toute interaction médicamenteuse ou dégradation du patient. 
Un dilemme persiste en réalité car dans un contexte palliatif terminal où Monsieur était hypotendu à l’admission avec une alimentation pauvre, est-ce que finalement le Dapagliflozine avait-il sa place ? Bien sûr la réponse reste à nuancer et l’initiation de ce traitement était clairement justifiable au niveau médical. 

\subsection{Plan de soins proposé} 

Deux alternatives étaient réalisables pour Monsieur. D’une part, une place au sein de l’unité de soins palliatifs la plus proche de son domicile était disponible. L’ensemble de son traitement pouvait être assuré en ces lieux. D’autre part, un retour à domicile pouvait être envisageable en considérant la mise en place d’aide à domicile et d’une collaboration étroite avec l’équipe mobile de soins palliatifs « Accompagner».

La décision prise entre l’équipe soignante, son épouse et Monsieur a été le retour à domicile. Parmi les aides mises en place, notons le rôle central du médecin traitant qui a su être à l’écoute et se rendre flexible afin de répondre au mieux aux demandes du couple. In fine, un passage quotidien de l’équipe infirmière, les transports ou le matériel médical éventuellement prodigué par la Croix-Rouge ou son épouse elle-même pour assurer la continuité des traitements, les aides ménagères ou alimentaires sont également proposées. Un protocole de détresse a pu être élaboré en cas de dégradation du patient. Malheureusement, Monsieur décédera 6 jours plus tard au domicile toutefois paisiblement grâce aux soins apportés, un suivi de deuil est également organisé pour son épouse.

\subsection{Rôle en tant que stagiaire médecin} 

J’ai pu agir concrètement en mettant en contact l’équipe sociale avec le patient et ses proches. J’ai eu l’occasion d’interagir directement avec l’équipe des soins continus et palliatifs pour mieux comprendre les enjeux et les options réalisables pour monsieur et son épouse.

Étant stagiaire, je disposais de plus de temps pour être au chevet du patient. Cela m’a permis une communication plus authentique avec monsieur. Nous avons abordé ses peurs. J’ai tenté d’appliquer plusieurs techniques communicationnelles afin de matérialiser ses douleurs, réaliser l’origine réelle de celles-ci et surtout identifier les ressources possibles. Souvent, plusieurs choses peuvent être mises en place sans coût ni excès de temps or cela réconforte grandement le patient. Une différence peut être faite par la façon d’aborder le patient, les informations transmises et les explications claires et ciblées aux demandes du patient.  

 
